\section{関連研究}

\subsection{日本語入力と省入力インタフェース}
日本語入力ではかな漢字変換が広く用いられ，読み入力から漢字かな交じり文への変換が研究されてきた．
既存の IME は変換精度の面では十分に成熟しており，残された課題はむしろ人間の側に要求される操作量にある．
モバイル環境ではフリック入力が普及し，小型端末向けの省入力手法も検討されている．
携帯端末向けの曖昧入力（T9 など複数文字を 1 キーに割り当てる方式）や，
予測に基づく加速入力（Dasher 等）は，少ない操作量での入力を目指す点で本研究と方向性を共有する．
視線入力や VR/AR 入力では，ポインティングや滞留選択の負荷が大きく，省入力の要請が強い．
これらの省入力インタフェースに共通するのは，端末や状況によらず人間の側に求められる操作を減らし，
より多くの人が負担少なく入力できるようにするという志向である．

\subsection{LLM と LoRA}
大規模言語モデルは文脈に基づく言語の補完・推定に高い性能を示す．
本研究は，曖昧な母音入力に対し，文脈を考慮した候補を生成する目的で LLM を用いる．
予測に基づく入力という観点では，曖昧入力の弁別や
言語モデル駆動の入力が先行するが，
本研究は子音を落とした母音列という極端な省入力を対象とする点が異なる．
モデルの追加学習には LoRA（Low-Rank Adaptation）を用いる．
LoRA は事前学習済み重みを固定し，低ランク行列のみを学習することで，
少ない学習パラメータで効率的にタスク適応を行う手法である．

\subsection{テキスト入力の評価}
テキスト入力手法は，打鍵効率の指標である KSPC（Keystrokes Per Character）や
誤り率により定量評価される．
主観的作業負荷の評価には NASA-TLX，
ユーザビリティの評価には SUS が広く用いられる．
本研究もこれらの枠組みに沿って評価設計を行う．
